\section{Rückblick}
\setauthor{Zoe Emily Öllinger}

Das Ziel unserer Diplomarbeit war ein Dashboard für die Firma Hello Again zu entwickeln, dass den Mitarbeitern hilft, die Health-Stati der Kundenapps zu überwachen.
Hierbei soll gewährleistet werden, dass die Mitarbeiter einen einfachen und intuitiven Workflow haben.
Daher wurden zwei Seiten erstellt.
Einmal die Übersichtsseite und einmal die Detailseite.


\section{Was wurde umgesetzt?}
\setauthor{Zoe Emily Öllinger}

Umgesetzt wurden eine Übersichtsseite, bei der man auf den ersten erkennen kann, ob eine App Healthy ist oder nicht, sowie die Anzahl der Tickets und wer den Account verwaltet.
Dabei wurden Quickfilter eingebaut mit mehrfachauswahl, damit die Mitarbeiter schnell nach ungesunden Apps filtern können.
Es gibt auch eine Suchleiste bei der man nach dem Slug filtern kann.

Auf der Detailseite, sieht man nochmal die genrellen Daten, die ebenfalls auf der Übersichtsseite angezeigt werden.
Danach gibt es die einzelnen Health-Checks, hier wird der Status der App und der Fehlercode angezeigt, sowie die dazugehörige Dokumentation um die Fehler zu beheben.
Es gibt eine Kommentarfunktion in der man Kommentare oder Tasks erstellen und zuweisen kann.
Weiters gibt einen Abschnitt in dem die vorhandenen Tickets angezeigt werden, sowie ein Banner, das im Falle einer Statusveränderung zum Vortag den neuen Status anzeigt.


\section{Wurden die Ziele erreicht?}
\setauthor{Zoe Emily Öllinger}

Unsere Ziele wurden im großen und ganzen erreicht.
Die Übersichtsseite und die Detailseite wurden so implementiert wie am Anfang geplant, wenn auch mit stylistischen Änderungen, da wir uns an das Styling des bereits vorhandenen Dashboards halten mussten.
Am Anfang wurde jedoch ein Analyse Tab geplant, der im laufe der Diplomarbeit allerdings von unseren Betreuern verworfen wurde.


\section{Herausforderungen und Schwierigkeiten}
\setauthor{Zoe Emily Öllinger}

Eine Herausforderung war es für uns, sich in einem Firmeninternen Environment zurechtzufinden und sich an die internen Codingguidelines zu halten, sowie mit neuen Programmiersprachen umzugehen, mit denen wir bis dorthin noch nicht in Berührung gekommen sind.
Für unsere Betreuer gab es auch Herausforderungen.
Am Anfang war es nicht ganz klar wo und wie wir arbeiten und dies zog sich über die ganze erste Woche des Praktikums.
Dabei gab es auch noch Hindernisse bezüglich den firmeninternen Accounts und Zugänge, die wir erst zu späterer Zeit erhalten haben.
Daher konnte in der ersten Woche auch noch nicht gleich zu 100 Prozent mit dem Coding gestartet werden.

\section{Was haben wir gelernt?}
\setauthor{Zoe Emily Öllinger}

Durch unsere Diplomarbeit, konnten wir viel praktisches Know-How lernen, da man gesehen hat, wie die Wirtschaft außerhalb der Schule wirklich funktioniert.
Außerdem lernten wir Developer kennen, die uns immer wieder Tipps gaben und auf Verbesserungen hinwiesen, die einem von selbst vielleicht garnicht aufgefallen wären.
Dazu haben wir durch die Workflow-Meetings ein gespür dafür bekommen, wie unterschiedlich, die Anforderungen von Kunden für eine Software sein können und wie man am besten alle Wünsche einbringt.
Wir durften uns außerdem auch noch mit vielen neuen Tools oder Frameworks beschäftigen wie zum Beispiel Vue.js und Asana.


\section{Könnte man unser Projekt in Zukunft noch erweitern?}
\setauthor{Zoe Emily Öllinger}

Man könnte unser Dashboard in Zukunft noch um weitere Fuktionen, wie zum Beispiel einen Statistik und Analyse Tab erweitern und weiter Tools einbauen.
Doch wie es mit dem Projekt weitergeht liegt nun in den Händen der Firma.